\paragraph{Communication sparsity in multi-agent reinforcement learning.}
Learned inter-agent communication has converged on attention for deciding whose
message matters: CommNet broadcasts mean-pooled hidden states
\citep{sukhbaatar2016commnet}, TarMAC introduces signature-based targeting
\citep{das2019tarmac}, and MAT casts the joint policy as a sequence model
\citep{wen2022mat}. Because attention over $N$ agents is pairwise, much of this
literature restricts who communicates, usually motivated by cost; DGN observes
that receiving much information ``requires high bandwidth and incurs high
computational complexity'' \citep{jiang2018dgn}. Corollary~\ref{cor:taxonomy}
sorts these mechanisms by what their selector costs. DGN's nearest neighbours,
ATOC's communication field and mean-field interaction fix the retained set from
structure available before any score exists
\citep{jiang2018dgn,jiang2018atoc,yang2018meanfield}; they are free if the
neighbourhood arrives with the environment and case~(b) if it is recomputed from
all pairwise positions each step. IC3Net and SchedNet gate or rank agents by a
scalar per agent, a top-$k$ over an $O(N)$ vector rather than an $N \times N$
matrix \citep{singh2018ic3net,kim2019schednet}. G2ANet and MAGIC run a pairwise
hard-attention or graph-construction network before communication attention
\citep{liu2020g2anet,niu2021magic}, so their selector is of the same order as
the scores they avoid. Score-then-select is the limiting case: the adaptive
sparse communication graph of \citet{sun2020sparsecomm} applies a sparsemax
generalisation to logits already formed for every pair and motivates it partly
by reduced complexity, and handshake mechanisms that match a query against every
candidate key before selecting \citep{liu2020who2com} are score-then-select in
the same sense. These methods operate with two to roughly thirty agents, DGN reaching the low
hundreds. Of the multi-agent papers we relied on that make an efficiency claim,
one of five reports a measured time attributable to the mechanism
(Appendix~\ref{app:audit}).

\paragraph{Efficient and sparse attention.} Efficient-attention methods
separate along the same axis. Sparse Transformers, Longformer and BigBird fix
the interaction pattern in advance, so the score computation itself is sparse
and the selector is essentially free
\citep{child2019sparse,beltagy2020longformer,zaheer2020bigbird}. Reformer and
Routing Transformer restrict candidates by hashing and online clustering, which
are content-dependent but evaluate scores only for candidate pairs, at a
selector cost that must be counted \citep{kitaev2020reformer,roy2021routing}.
Performer keeps every interaction and reorders the algebra
\citep{choromanski2021performers}. Top-$k$ attention in sequence models is the
closest analogue of our mechanism: the Explicit Sparse Transformer keeps the $k$
largest scores per query \citep{zhao2019explicit}, and \citet{gupta2021topk}
compute each query's exact top-$k$ over all keys in chunks, with memory
efficiency as the stated benefit. Sparsemax and $\alpha$-entmax obtain exact
zeros from the normalisation rather than a hard cutoff
\citep{martins2016sparsemax,peters2019entmax}; because their zeros arrive after
the scores exist, they fall under Proposition~\ref{prop:ceiling}, which is our
inference rather than a claim they make. Sparse mixture-of-experts routing is
also score-then-select \citep{shazeer2017moe,fedus2022switch}, but the router is
small relative to the expert computation it gates, so the analogue of our score
term is minor and the saving is real. Whether score-then-select pays therefore
depends on the ratio between the cost of scoring every candidate and the cost
of the work skipped, which is one in single-head attention. 

\paragraph{Hardware-aware execution.} That operation counts do not determine
runtime is established in general terms \citep{dehghani2022misnomer}, and
transformer workloads are known to be dominated by data movement rather than
arithmetic \citep{ivanov2021datamovement}; the roofline model formalises when
arithmetic ceases to govern runtime \citep{williams2009roofline}. FlashAttention
and its successors reduce memory traffic without changing the operation count
\citep{dao2022flashattention,dao2024flashattention2,shah2024flashattention3,rabe2021selfattention},
and FlexAttention compiles attention variants expressed as score modifications
and block masks into fused kernels \citep{dong2025flexattention}. Native Sparse
Attention and MiniMax Sparse Attention show where sparse attention does pay:
blockwise selection amortises the selector over blocks, preserves matrix-multiply
shapes inside selected tiles, and operates at $10^4$ to $10^6$ tokens, where the
skipped work is large \citep{yuan2025nsa,lai2026msa}. Relative to this
literature our contribution is specific: a selector-cost accounting that
separates mechanisms before measurement, a byte criterion for per-pair
compaction, and measurements at tens to hundreds of agents and narrow heads,
where neither block amortisation nor large skipped work is available. Our benchmarking follows
established practice for reporting performance on shared hardware
\citep{hoefler2015benchmarking}, and our kernels are written in Triton
\citep{tillet2019triton}. The training study follows
\citet{papoudakis2021epymarl} and \citet{agarwal2021rliable}, with environments
from PettingZoo/MPE \citep{terry2020pettingzoo} and MAPPO
\citep{yu2022mappo}.
